\section{Experimental design}
\label{sec:experiments}
\subsection{Development protocol and landmark evaluation}
The 801 development patients were partitioned into five patient-level folds. Fold 1 used 640 training and 161 validation patients; folds 2--5 each used 641 training and 160 validation patients. The same 175-patient holdout was evaluated by every fold model and every variant. Thus the five test evaluations per variant are repeated evaluations of overlapping trained models on the same patients, rather than five independent test cohorts. Patient identities were disjoint between training, validation, and holdout within the current split.

At each training epoch, a patient contributed one randomly selected available T0, T1, or T2 prefix. Only MRI observations available at or before the selected landmark were classifier inputs. Later observed visits, including T3, could supply latent targets for the auxiliary loss on training patients. Missing targets were masked. The 23 patients with only T0 remained eligible for pCR training but supplied no future-state loss. At evaluation, all three landmarks were evaluated for the 175 holdout patients, each of whom had T0, T1, and T2. T3 was not an evaluation landmark and was not used as observed input for these predictions.

\subsection{Model comparisons}
The principal comparison comprises three variants trained with seed44 and the same five development folds. \nojepa{} uses direct pCR readout and a zero future-state loss weight. \direct{} uses the same direct readout with a future-state loss weight of 0.5. \fullmodel{} retains that auxiliary weight and replaces the direct readout with a concatenation of direct pCR features, attention-pooled predicted terminal state, and attention-pooled predicted change. All variants retain the shared image encoder, clinical inputs, assigned-arm encoding, fusion architecture, and training schedule.

The \direct{}--\nojepa{} contrast most directly tests adding future-state supervision. The \fullmodel{}--\direct{} contrast jointly changes the use of forecast features, spatial pooling, and classifier input width. It cannot isolate attention alone. For that purpose, the supplementary \meanmodel{} comparison uses forecast readout with mean spatial pooling and otherwise matches the forecast-attention configuration. All reported models in this draft use seed44. No VQ-latent reconstruction-head, change-weighted-loss, or state-conditioned-head variants are included.

\subsection{Implementation and optimization}
The network uses a model width of 768, eight layers in each of the spatial encoder, multimodal fusion module, and predictor, 16 attention heads with 64 dimensions per head, and feed-forward width 3,072. The 1,024-dimensional concatenated attention-head representation is projected back to the 768-dimensional residual stream. The spatial token grid is $2\times4\times4$. The configuration uses eight state queries, four pCR queries, dropout 0.3, and an EMA coefficient of 0.999. The pretrained convolutional encoder is fine-tuned jointly with the remaining network. Its training-patient provenance requires additional documentation, as discussed in Section~\ref{sec:limitations}.

We use AdamW \citep{adamw2019} with an initial base learning rate of $10^{-4}$ for both the encoder and the main network, weight decay 0.05, batch size 16, no gradient accumulation, and gradient-norm clipping at 1.0. Training uses BF16 mixed precision. A ten-epoch warm-up increases the learning rate from $10^{-5}$ to $10^{-4}$, followed by cosine decay to $10^{-6}$ at epoch 100. Training is capped at 200 epochs, with early stopping after 50 epochs without improvement. The best checkpoint is selected independently in each fold using validation patient-equal-prefix F1 at threshold 0.5. This is an F1 selection protocol, even though AUROC is a principal reporting measure. Deterministic execution was not enabled, so a fixed seed does not imply bitwise reproducibility across hardware or software changes.

\subsection{Metrics, aggregation, and patient-level uncertainty}
We report AUROC, average precision (AP), and F1 at threshold 0.5 at each landmark. AP is calculated as the non-interpolated average precision statistic. To assess probability quality, we also report negative log-likelihood (NLL), the Brier score, and expected calibration error (ECE) with ten equal-width probability bins. No test-set threshold optimization or probability recalibration was performed. The positive prevalence of the holdout is $56/175=0.32$.

Let $M$ denote a metric, $p_{if\ell}$ the probability for patient $i$ from fold model $f$ at landmark $\ell$, and $y_i$ the pCR label. Our landmark summary is
\begin{equation}
 \overline M_{\ell}=\frac{1}{5}\sum_{f=1}^{5}M\bigl(\{y_i,p_{if\ell}\}_{i=1}^{175}\bigr),
 \qquad
 \overline M_{\mathrm{macro}}=\frac{1}{3}\sum_{\ell=0}^{2}\overline M_{\ell}.
\end{equation}
We do not average fold probabilities before computing $M$. Five-fold metric means and probability-ensemble metrics are distinct estimands. Similarly, a mean F1 across landmarks is not recomputed after pooling all landmark rows.

For exploratory uncertainty estimates, we draw 2,000 bootstrap samples of the 175 patients with replacement. Each patient draw is shared across all fold models, variants, and landmarks, preserving pairing and repeated observations. Within each resample, metrics are computed separately per fold and then averaged exactly as above. Percentile 2.5th and 97.5th quantiles form the reported intervals. These intervals condition on the trained seed44 models and the selected analysis; they do not capture uncertainty from retraining, choosing a seed, or selecting a model after results were inspected. We do not use fold-level or landmark-level values as independent patient samples, and we do not make multiplicity-adjusted significance claims.

\subsection{Scope of the selected-seed analysis}
The choice of seed44 for this manuscript was made retrospectively after results from multiple training seeds were available. It was not a prospectively registered seed or a selection made without knowledge of the holdout results. Reporting the same seed for all models preserves a matched comparison within that seed, but it does not remove selection bias or establish robustness across training randomness. The analysis therefore evaluates an exploratory model configuration rather than providing a confirmatory performance estimate. A final study should lock model selection before evaluating an independent cohort and report a prespecified multi-seed sensitivity analysis.
